\section{Seguimiento Vehicular y Determinación de Eventos}

En un sistema de control de acceso para parqueaderos, la detección frame a frame
de los vehículos resulta insuficiente si no se acopla a un mecanismo de
continuidad temporal. El módulo de seguimiento vehicular (\textit{Vehicle
Tracking}) tiene como objetivo vincular las detecciones espaciales a lo largo
del tiempo para construir la trayectoria continua de cada automóvil. Esta etapa
es fundamental para transformar coordenadas crudas de píxeles en eventos
operativos estructurados, tales como la autorización de paso, el conteo en
tiempo real y la clasificación direccional de ingresos y salidas.

\subsection{Formulación matemática del seguimiento de objetos múltiples (MOT)}

Como sintetizan Luo et al. \cite{luo2021mot}, el seguimiento de objetos
múltiples (\textit{Multiple Object Tracking}, MOT) se formula como un problema
de estimación multivariable donde, dada una secuencia continua de video, se
busca determinar el estado de $M_t$ objetos presentes en cada fotograma $t$. En
el paradigma estándar de seguimiento por detección (\textit{Tracking-by-Detection}), 
el detector de objetos (como YOLO) proporciona en cada frame $t$ un conjunto de
observaciones $O_t = \{o_1^t, o_2^t, \dots, o_{N_t}^t\}$, mientras que el sistema
de seguimiento mantiene el conjunto de estados hipotéticos $S_t = \{s_1^t, s_2^t, \dots, s_{M_t}^t\}$.

La tarea central del problema MOT reside en la asignación temporal de identidades:
asignar y mantener un identificador único e invariable ($ID_k \in \mathbb{N}$) a
cada vehículo individual a medida que se desplaza dentro del campo de visión de
la cámara. 

En la operación real de un parqueadero, la asignación de este $ID$ único es
indispensable. Si el sistema dependiera únicamente de detecciones independientes
por frame a una tasa de 30 fotogramas por segundo, un automóvil detenido o
avanzando lentamente hacia la barrera generaría decenas de registros aislados
por segundo. Mantener una trayectoria persistente (\textit{tracklet}) mediante
un $ID$ único garantiza que el ciclo de vida del vehículo sea tratado como una
sola entidad continua desde que entra en el área de aproximación hasta que cruza
la barrera, evitando la duplicación estocástica de eventos de acceso.

\subsection{Modelado de incertidumbre y estimación del estado mediante Filtro de Kalman}

Como demostró originalmente Kalman \cite{kalman1960new} y adaptaron Bewley et al.
\cite{bewley2016sort} para el seguimiento en tiempo real, el movimiento de una
caja delimitadora puede modelarse dinámicamente mediante un Filtro de Kalman
lineal bajo la suposición de velocidad constante. Para el caso de un vehículo,
el vector de estado sintético $s_t$ se define típicamente en el espacio de la
imagen como:

\begin{equation}
s_t = [u, v, s, r, \dot{u}, \dot{v}, \dot{s}]^T
\end{equation}

donde $(u, v)$ representan el centroide espacial de la caja delimitadora en
píxeles, $s$ es la escala o área total de la caja ($s = w \times h$), $r$ es la
relación de aspecto fija ($r = w / h$), y los términos $(\dot{u}, \dot{v}, \dot{s})$
representan sus respectivas velocidades lineales de cambio entre fotogramas
consecutivos.

El Filtro de Kalman opera de manera recursiva en dos fases conceptuales continuas:

\begin{enumerate}
    \item \textbf{Fase de Predicción (\textit{Predict}):} Utilizando la ecuación
        del modelo cinemático previa, el filtro proyecta hacia adelante el
        estado estimado y su matriz de covariaza de incertidumbre para el
        fotograma $t+1$, anticipando dónde debería ubicarse el vehículo antes de
        procesar la nueva imagen.
    \item \textbf{Fase de Actualización (\textit{Update}):} Una vez que el
        detector emite las observaciones reales del fotograma $t+1$, el filtro
        calcula la ganancia de Kalman para ponderar la incertidumbre del sensor
        frente a la predicción teórica, ajustando corregidamente el estado final
        del vehículo.
\end{enumerate}

Esta capacidad de estimación cinemática proporciona al sistema un mecanismo de
``memoria temporal''. En el acceso a un parqueadero, es habitual que la placa o
el parachoques del vehículo sufran oclusiones parciales temporales —por ejemplo,
cuando un peatón camina frente a la cámara, cuando la barrera física se levanta
o ante destellos de luz solar directa—. Durante los fotogramas en los que el
detector no genera ninguna observación ($o_t = \emptyset$), el Filtro de Kalman
mantiene la predicción de la posición del vehículo en función de su velocidad
previa, evitando que la trayectoria se rompa o que el sistema pierda la
identidad del automóvil durante breves intervalos de ceguera óptica.

\subsection{Asociación de datos y optimización combinatoria: Algoritmo Húngaro y SORT}

Para enlazar las predicciones proyectadas por el Filtro de Kalman con las nuevas
detecciones generadas por el modelo de visión en cada frame, Bewley et al.
\cite{bewley2016sort} formularon el algoritmo SORT (\textit{Simple Online and
Realtime Tracking}). SORT aborda la asociación de datos resolviendo un problema
de asignación lineal global mediante el Algoritmo Húngaro (Kuhn)
\cite{kuhn1955hungarian}, una técnica de
optimización combinatoria que minimiza el costo total de emparejar dos conjuntos
disjuntos.

Dadas $M$ trayectorias predichas por Kalman ($\hat{B}_{\text{pred}}$) y $N$
nuevas cajas detectadas por YOLO ($B_{\text{det}}$), se construye una matriz de
costo asociativo $C \in \mathbb{R}^{M \times N}$. La métrica elegida para medir
la distancia espacial entre la predicción $i$ y la detección $j$ es el
complemento de la Intersección sobre Unión (IoU):

\begin{equation}
C_{i,j} = 1 - \text{IoU}\left(\hat{B}_{\text{pred}, i}, \; B_{\text{det}, j}\right)
\end{equation}

El Algoritmo Húngaro procesa la matriz $C$ para encontrar la asignación óptima
que minimice la suma global de distancias espaciales. Para evitar
emparejamientos erróneos entre vehículos lejanos, se aplica un umbral rígido de
rechazo ($\text{IoU}_{\text{min}}$, típicamente $0{.}3$):

\begin{itemize}
    \item Si $C_{i,j} > 1 - \text{IoU}_{\text{min}}$, la asignación se descarta.
    \item Las detecciones no asociadas a ninguna trayectoria existente se
        inicializan como candidatos a nuevos vehículos.
    \item Las trayectorias que no reciben ninguna detección asociada durante un
        número límite de fotogramas ($T_{\text{lost}}$) son eliminadas
        definitivamente de la memoria del sistema.
\end{itemize}

Esta combinación de baja complejidad computacional permite a SORT resolver la
asociación de datos en una fracción de milisegundo por frame, manteniendo el
procesamiento en tiempo real.

\subsection{Determinación de eventos a partir de trayectorias vehiculares}

La etapa final del pipeline transforma las trayectorias asociadas ($Tracklets$)
en eventos discretos con validez operativa para la gestión del parqueadero. Como
analizan Liu y Zeng \cite{liu2017vehiclecounting} y Karbalaie et al.
\cite{karbalaie2022event}, la lógica
de negocio para la detección de eventos se fundamenta en la definición de líneas
virtuales de control en el plano de la imagen, conocidas técnicamente como
\textit{Virtual Tripwires}.

Una línea virtual $L_{\text{tripwire}}$ se define geométricamente como un
segmento orientado en el espacio de la imagen mediante dos puntos de control
$P_1(x_1, y_1)$ y $P_2(x_2, y_2)$, ubicada estratégicamente sobre la calzada
frente a la barrera de acceso.

A medida que un vehículo con identidad $ID_k$ avanza, el sistema registra la
secuencia histórica de los centroides de su caja delimitadora $\mathbf{p}_t =
(u_t, v_t)$ para instantes $t_1 < t_2 < \dots < t_n$. La determinación del
evento se evalúa geométricamente mediante el cruce de línea:

\begin{enumerate}
    \item \textbf{Detección del Cruce Espacial:} Se verifica analíticamente si
        el segmento trazado entre la posición previa del vehículo $\mathbf{p}_{t-1}$
        y su posición actual $\mathbf{p}_t$ se interseca espacialmente con el
        segmento de la línea virtual $L_{\text{tripwire}}$.
    \item \textbf{Cálculo del Vector de Dirección:} Una vez confirmado el cruce,
        se evalúa el producto escalar o el producto cruz vectorial entre el
        vector de desplazamiento del vehículo $\mathbf{v}_{\text{mov}} =
        \mathbf{p}_t - \mathbf{p}_{t-1}$ y el vector normal unitario de la línea
        virtual $\mathbf{n}_L$:
    
    \begin{equation}
    \text{Dirección} = \text{signo}\left( \mathbf{v}_{\text{mov}} \cdot \mathbf{n}_L \right)
    \end{equation}
\end{enumerate}

Dependiendo del signo del resultado geométrico, el sistema clasifica el evento
de manera determinista:

\begin{itemize}
    \item \textbf{Evento de Ingreso (\textit{Entry Event}):} Si el vector de
        movimiento apunta hacia el interior del parqueadero
        ($\text{Dirección} > 0$), el sistema consolida el evento registrando el
        $ID$ del vehículo, la lectura de la placa vinculada, la fecha, la hora
        exacta y habilita la apertura de la barrera.
    \item \textbf{Evento de Salida (\textit{Exit Event}):} Si el vector de
        movimiento apunta en sentido opuesto ($\text{Dirección} < 0$), el
        sistema registra la salida, calcula el tiempo de permanencia consumido
        en el recinto y procede al cobro o liberación de la plaza.
\end{itemize}

Mediante esta arquitectura integrada de detección, filtrado cinemático,
asociación combinatoria y análisis geométrico de trayectorias, el sistema
garantiza un registro de accesos automatizado, auditable y resistente a
falsas disparaciones.
